\documentclass[preview]{standalone}

\usepackage{amsmath}
\usepackage{amssymb}
\usepackage{stellar}
\usepackage{definitions}

\begin{document}

\id{jordan-block-powers-ranks}
\genpage

\section{Nilpotent operators}

\begin{snippetdefinition}{nilpotent-endomorphism-definition}{Nilpotent endomorphisms and matrices}
    A \linearendomorphism \(N\colon V\fromto V\) is
    \emph{nilpotent} if
    \[
        N^r=0
    \]
    for some integer \(r\geq1\). The least such \(r\) is its
    \emph{nilpotency index}. A square \matrix is nilpotent if the same
    condition holds for one of its positive powers, equivalently if its
    associated endomorphism is nilpotent.
\end{snippetdefinition}

\begin{snippetproposition}{jordan-block-nilpotent-part}{Nilpotent part of a Jordan block}
    A \jordanblock can be written as
    \[
        J_m(a)=aI_m+N_m,
    \]
    where
    \[
        N_m
        =
        \begin{pmatrix}
            0&1&0&\cdots&0\\
            0&0&1&\cdots&0\\
            0&0&0&\ddots&\vdots\\
            \vdots&\vdots&\ddots&\ddots&1\\
            0&0&\cdots&0&0
        \end{pmatrix}.
    \]
    The matrix \(N_m\) is \nilpotent of index \(m\):
    \[
        N_m^m=0,
        \qquad
        N_m^{m-1}\neq0.
    \]
\end{snippetproposition}

\begin{snippetproof}{jordan-block-nilpotent-part-proof}{jordan-block-nilpotent-part}{Nilpotent part of a Jordan block}
    On the standard basis,
    \[
        N_me_1=0,
        \qquad
        N_me_i=e_{i-1}
        \quad(2\leq i\leq m).
    \]
    Hence \(N_m^ke_i=e_{i-k}\) when \(i>k\), and it is zero when
    \(i\leq k\). In particular, \(N_m^m=0\), while
    \(N_m^{m-1}e_m=e_1\neq0\).
\end{snippetproof}

\section{Ranks of powers}

\begin{snippetproposition}{nilpotent-jordan-block-power-rank}{Ranks of powers of a nilpotent Jordan block}
    For every \(k\in\naturalnumbers\),
    \[
        \mrank(N_m^k)
        =
        \begin{cases}
            m-k,&0\leq k\leq m,\\
            0,&k\geq m,
        \end{cases}
    \]
    and equivalently
    \[
        \dim_K\ker(N_m^k)=\min\{k,m\}.
    \]
\end{snippetproposition}

\begin{snippetproof}{nilpotent-jordan-block-power-rank-proof}{nilpotent-jordan-block-power-rank}{Ranks of powers of a nilpotent Jordan block}
    The formula
    \[
        N_m^ke_i
        =
        \begin{cases}
            e_{i-k},&i>k,\\
            0,&i\leq k
        \end{cases}
    \]
    shows that, for \(k\leq m\), the image is generated by
    \(e_1,\ldots,e_{m-k}\). For \(k\geq m\), the image is zero. The
    kernel formula follows from rank--nullity.
\end{snippetproof}

\begin{snippettheorem}{jordan-block-sizes-from-kernel-dimensions}{Recovering Jordan blocks from kernel dimensions}
    Let \(A\in\matrices_n(K)\), suppose its characteristic polynomial
    splits over \(K\), and let \(a\in K\) be an \eigenvalue. Put
    \[
        \nu_k=\dim_K\ker((A-aI_n)^k),
        \qquad
        \nu_0=0.
    \]
    Then
    \[
        \nu_k-\nu_{k-1}
    \]
    is the number of Jordan blocks for \(a\) having size at least \(k\).
    Consequently,
    \[
        2\nu_k-\nu_{k-1}-\nu_{k+1}
    \]
    is the number of blocks for \(a\) having size exactly \(k\).
    Thus the sequence \(\nu_1,\nu_2,\ldots\) determines all Jordan-block
    sizes for \(a\).
\end{snippettheorem}

\begin{snippetproof}{jordan-block-sizes-from-kernel-dimensions-proof}{jordan-block-sizes-from-kernel-dimensions}{Recovering Jordan blocks from kernel dimensions}
    Similarity preserves the dimensions of these kernels, so compute in
    Jordan form. On a block \(J_m(a)\), the operator \(A-aI_n\) becomes
    \(N_m\), and its contribution to \(\nu_k\) is \(\min\{k,m\}\).
    Therefore its contribution to
    \(\nu_k-\nu_{k-1}\) is \(1\) exactly when \(m\geq k\).

    On a block \(J_m(b)\) with \(b\neq a\), the matrix
    \[
        J_m(b)-aI_m=(b-a)I_m+N_m
    \]
    is invertible, so it contributes nothing to the kernel. Summing
    over the \(a\)-blocks proves the first formula. Subtracting the
    number of blocks of size at least \(k+1\) from the number of blocks
    of size at least \(k\) gives the exact-size formula.
\end{snippetproof}

\section{Example}

\begin{snippetexample}{jordan-block-ranks-two-eigenvalues-example}{Reading blocks from ranks}
    Suppose \(a\neq b\) and
    \[
        A\sim
        J_1(a)\oplus J_1(a)\oplus J_3(a)
        \oplus J_2(b)\oplus J_2(b).
    \]
    The total dimension is \(9\). On the \(a\)-primary component, the
    ranks of the first three powers of \(A-aI\) are
    \(2,1,0\); on the \(b\)-primary component, \(A-aI\) is invertible
    and contributes rank \(4\). Hence
    \[
        \mrank(A-aI)=6,
        \qquad
        \mrank((A-aI)^2)=5,
        \qquad
        \mrank((A-aI)^3)=4.
    \]
    The corresponding nullities are \(3,4,5\), whose successive
    differences are \(3,1,1\). Thus there are three \(a\)-blocks, one
    of size at least \(2\), and one of size at least \(3\); their sizes
    are
    \[
        1,\ 1,\ 3.
    \]

    Similarly,
    \[
        \mrank(A-bI)=7,
        \qquad
        \mrank((A-bI)^2)=5,
        \qquad
        \mrank((A-bI)^3)=5.
    \]
    The nullities are \(2,4,4\), so there are two \(b\)-blocks and both
    have size at least \(2\). Their sizes are \(2,2\).
\end{snippetexample}

\end{document}
